% Transcription — fonds Alexandre Grothendieck, shelfmark 103, batch 1
% (pages 1-17, the whole folder). Established from the facsimile
% archives/batches/103/batch-01.pdf (a mirror of
% https://grothendieck.umontpellier.fr/103.pdf), per the skill
% transcribe-grothendieck. The batch file carries no cover sheet: PDF page
% and archivists' page coincide.
%
% The folder holds three documents, and no working notes.
%
% Pages 1 to 6: a six-page typescript, « Esquisse d'une théorie des
% Gr-Catégories, par Mme Hoang Xuan Sinh » — a résumé of Hoàng Xuân Sính's
% thesis (whose typescript is folder 135), corrected and annotated by hand:
% interlinear insertions (« c'est un groupoïde, et si », « ponctué par
% x ∈ X », « localement », « plus généralement », « axiome du pentagone »,
% « non canonique », « Deligne » filling the empty citation brackets),
% struck typed words, a marginal note on page 1 and a struck paragraph on
% page 5. The annotating hand is, to all appearance, the one that writes the
% letter of pages 8-15; the page does not sign it and the file does not
% assert it beyond that. The top right of page 1 carries, in another and
% upright hand, « 9/74 Envoyé par Grothendieck » — a recipient's docket.
% Following folders 48 and 55, an annotated typescript gives the statements
% it carries and his interventions, and summarises rather than recomposes
% typed proofs. This typescript is itself a résumé — definitions, the
% structure theorem, remarks, questions — and carries no proof; there is
% nothing to summarise away, and it is given whole. His insertions are set
% \add{} and the struck typed words \struck{}, as in folder 55, with a note
% at page 1 saying so; typed underlinings are \emph{}.
%
% Page 7 is blank and is skipped; the gap in the numbering is the record.
%
% Pages 8 to 15: a letter in Grothendieck's hand, « Lodève le 27.8.74 »,
% « Cher Deligne », eight pages. Its first four pages (8-11) are the letter
% folder 134-1 holds as pages 3-6 — the copy Deligne forwarded to Breen; this
% copy runs four pages further. Pages 8-11: the séminaire d'algèbre on « les
% fourbis de Mme Sinh » transposed to champs; E(M,N) = tau_{<=2} RHom(M,N),
% the triangle (T) with E'(M,N) and Hom(M, 2N)[-2], the exact sequence (*)
% 0 -> Ext^2 -> P(M,N) -> Hom(M, 2N) -> 0 and its heuristic reading as the
% strict Picard 2-champ of not-necessarily-strict Picard 1-champs pinned by
% M, N; the example of K^0, K^1 of the additive champ of projective
% A-Modules; the question whether (T) and (*) are known to Quillen, Breen,
% Illusie. Pages 14-15: Gr-champs on X pinned by (G, N), their 2-category
% expressed by tau_{<=2}(RGamma(B_G mod X, N)[1]) and read as 2-gerbes on B_G
% bound by N trivialised over X = B_e = (B_G)_{/P}; the localised complex
% tau_{<=2}(R p_{G*} Coker(N -> R q_{G*} C(q_G^* N))). Pages 12-13: the
% question of relative cohomology RGamma(Y mod X, F) for a morphism of topoi
% q : X -> Y, with hand constructions in the two extreme cases — q
% (-1)-acyclic, q an open immersion — and the obstacle that the cone is not
% functorial in the derived category; a « Question pour Illusie » on
% deformations of flat group schemes via Gr-champs or Picard champs.
%
% The archivists' order of pages 12-13 and 14-15 is probably not the order
% of writing: page 12 refers to « B_e -> B_G plus haut », and B_G appears
% only on pages 14-15; page 13 ends on « J'ai » and page 14 opens « Je te
% signale », so the run 8-11, 14-15, 12-13 reads continuously and the leaf
% that followed 13 is not in the folder. Transcribed in the archivists'
% order, flagged where it occurs, not repaired.
%
% The date at the head of page 8 is faint: « 7.8.74 » is dark, and before
% the 7 there is only a short pale stroke, the foot of a 2 whose upper part
% has faded. Folder 134-1 is another reproduction of the same leaves, and on
% it the 2 is plain: « Lodève le 27.8.74 ». The first pass read « 7.8.74 »;
% on 2026-09-14 the date was revised to 27.8.74 by Opus 5 (claude-opus-5),
% the 2 marked uncertain since this copy alone does not show it, and the
% notes at pages 8 and 10 that reported 134-1's first-pass readings were
% brought up to date with that file's revision of 2026-09-13. Running up the left margin of page 8, in a hand not
% distinguished from the letter's: « Je n'ai pas retrouvé la lettre sur les
% Gr-champs que tu cites » — transcribed as marginal, with its uncertainty.
%
% Pages 16-17: a typed letter, « Villecun le 23.6.1974 », « Chers Deligne,
% Giraud, Verdier », signed « Schurik » by hand, asking them to sit on the
% jury of Mme Sinh's thesis. It is his letter, and following folder 134-1
% (his letter to Deligne given whole, séminaire and cours de 1er cycle
% included) it is transcribed whole, since its one paragraph of mathematics
% — the thesis subject as he states it — sits inside the rest. One line is
% withheld: the postal address and telephone number of Mme Sinh's brother,
% typed on page 17, which a private person's data has no place in an
% edition; a note marks the omission. No letter by another hand is in the
% folder, so the question of letters by other people does not arise here.
%
% Notation fixed once for the batch. His script H for hypercohomology is
% set \mathbb{H}; the typed « ⊗ » (an overstruck O on the machine) is
% \otimes; the underlined O of the structure sheaf is \underline{O}_X, as in
% folder 55; the 2-torsion he writes with a prefixed index is {}_2 N,
% {}_2(\pi_1). The typescript writes « chaines », « interpétation »,
% « Postnikoff », « Mac-Lane », « Picards » and « la "champ" »: left as
% typed.
%
% The dating is the inventory's own, for the folder.
%
% Pass: Fable 5.1 (claude-fable-5-1), 2026-09-13 — first pass, unchecked against the pages by a human.
\documentclass[11pt,a4paper]{article}
% Le préambule commun à toutes les transcriptions du fonds Grothendieck.
%
% Il tient en une idée : **l'incertitude se compose, elle ne se résout pas.**
% Une transcription qui lisse un mot douteux a détruit la seule chose pour
% laquelle on la faisait — un lecteur doit pouvoir savoir, à chaque endroit,
% ce qui était sur le feuillet et ce qui a été ajouté. Les six macros qui
% suivent sont donc l'appareil critique, et non de la décoration.
%
% Les mêmes commandes sont reconnues par `scripts/render.mjs`, qui produit la
% vue de lecture du site. Une macro ajoutée ici doit l'être aussi là-bas,
% faute de quoi le rendu échouera bruyamment — ce qui est voulu : un rendu
% silencieusement faux est pire qu'un rendu absent.

\ProvidesPackage{grothendieck}[2026/08/08 Transcriptions du fonds Grothendieck]

\usepackage{amsmath,amssymb,amsthm}
\usepackage{mathrsfs}
\usepackage{stmaryrd}
% Les diagrammes commutatifs sont partout dans ce fonds : c'est la langue
% même de la géométrie algébrique de Grothendieck, et une transcription qui
% les rendrait en prose ne transcrirait rien.
\usepackage{tikz-cd}
% Français par défaut — c'est la langue des feuillets — mais l'anglais est
% chargé aussi : la traduction et le résumé sont des documents anglais, et la
% typographie française (espaces avant les deux-points) y serait une faute.
% Ces documents basculent par \selectlanguage{english} en tête de corps.
\usepackage[english,french]{babel}
\usepackage{fontspec}
\usepackage{geometry}
\geometry{a4paper,margin=25mm,marginparwidth=28mm}
\usepackage{marginnote}
\usepackage{xcolor}
% ulem, not soul. `soul`'s \st and \ul are built on a character-by-character
% scan that XeLaTeX's font machinery breaks: the document fails with an
% undefined control sequence rather than degrading. ulem does the same two jobs
% and is Unicode-engine safe. `normalem` stops it hijacking \emph, which the
% transcriptions use for Grothendieck's own emphasis.
\usepackage[normalem]{ulem}
\usepackage{hyperref}
\hypersetup{colorlinks=true,linkcolor=black,urlcolor=black,pdfborder={0 0 0}}

% --- Métadonnées du lot -----------------------------------------------------
% Reprises telles quelles par le script de rendu, qui les affiche en tête de
% la vue de lecture. Elles fixent ce que le fichier prétend couvrir : sans
% elles, une transcription flotte sans point d'attache dans un fonds de seize
% mille feuillets.

\newcommand{\folder}[1]{\gdef\grothFolder{#1}}
\newcommand{\batch}[1]{\gdef\grothBatch{#1}}
\newcommand{\leaves}[2]{\gdef\grothFirst{#1}\gdef\grothLast{#2}}
\newcommand{\foldertitle}[1]{\gdef\grothFoldertitle{#1}}
\gdef\grothFolder{?}\gdef\grothBatch{?}\gdef\grothFirst{?}\gdef\grothLast{?}\gdef\grothFoldertitle{}

% --- Le résumé --------------------------------------------------------------
%
% Il ouvre la lecture modernisée, sous le titre « Résumé », et il est composé
% en petit. Ce n'est pas un ornement : c'est le seul passage du document qui
% ne soit pas des mathématiques, et un lecteur doit pouvoir le reconnaître
% avant de l'avoir lu — puis le sauter s'il connaît déjà le sujet, ou s'y
% arrêter s'il ne le connaît pas. Le filet qui le suit marque la reprise du
% corps.

\newenvironment{resume}
  {\small\setlength{\parskip}{0.45em}}
  {\par\medskip\hrule\medskip}

% --- Mots-clés ---------------------------------------------------------------
%
% En anglais, en clôture du résumé de la lecture modernisée : le vocabulaire
% moderne sous lequel chercher ce que les feuillets construisent. Le manifeste
% du site (`npm run manifest`) les extrait pour étiqueter la cote entière —
% cette ligne est la source unique des tags, il n'y a pas de fichier à côté.
\newcommand{\keywords}[1]{\par\smallskip\noindent
  {\footnotesize\textsc{Keywords} --- \textit{#1}}\par}

% --- L'appareil critique ----------------------------------------------------

% Le numéro de feuillet, en marge. C'est l'ancre qui permet de régler un
% désaccord sur une formule en regardant *un* feuillet plutôt qu'une chemise
% de six cents. Le site s'en sert aussi pour tourner les pages du fac-similé
% quand on fait défiler la transcription.
\newcommand{\leaf}[1]{%
  \marginnote{\footnotesize\color{gray!70}\textsf{#1}}%
  \label{leaf:#1}\ignorespaces}

% Illisible : on ne devine pas. Un mot inventé qui se lit comme les autres est
% le pire résultat possible.
\newcommand{\ill}{\textcolor{red!60!black}{[\,\ldots\,]}}

% Lecture douteuse : le mot est donné, et signalé comme incertain.
\newcommand{\uncertain}[1]{\uline{#1}}

% Ajout de l'éditeur — une lettre manquante, un mot sous-entendu.
\newcommand{\add}[1]{\textcolor{blue!50!black}{[#1]}}

% Biffé par Grothendieck lui-même. Ce qu'il a rayé fait partie du document :
% une rature dit souvent où la pensée a changé de direction.
\newcommand{\struck}[1]{\sout{#1}}

% Note du transcripteur, sur l'état matériel du feuillet.
\newcommand{\note}[1]{\par\smallskip{\small\color{gray!60!black}\textit{[#1]}}\par\smallskip}

% Note marginale de Grothendieck, distincte du corps du texte.
\newcommand{\marginal}[1]{\par\smallskip\noindent\hspace{1em}%
  {\small\color{gray!60!black}#1}\par\smallskip}

% --- Titre ------------------------------------------------------------------

\AtBeginDocument{%
  \begin{center}
    {\large\bfseries Fonds Alexandre Grothendieck --- cote n\textsuperscript{o}~\grothFolder}\\[2pt]
    {\itshape\grothFoldertitle}\\[4pt]
    {\small feuillets \grothFirst--\grothLast\ (lot \grothBatch)}\\[2pt]
    {\footnotesize\color{gray!60!black}Transcription établie d'après le fac-similé numérisé
     par l'Université de Montpellier.\\
     Les lectures incertaines sont soulignées, les illisibles notées [\,\ldots\,],
     les ajouts entre crochets.}
  \end{center}
  \vspace{1em}}

\endinput


\folder{103}
\batch{1}
\pages{1}{17}
\dating{1974}
\watermark{Édition de démonstration}
\foldertitle{Courrier [Correspondance sur la thèse de Sinh] : copies d'article annoté (1974), lettres (1974)}

\begin{document}

\section*{Esquisse d'une théorie des Gr-Catégories, par Mme Hoang Xuan Sinh --- tapuscrit annoté (pages 1 à 6)}

\page{1}
\note{Tapuscrit de six pages, corrigé et annoté à la main. Comme pour les
tapuscrits des dossiers 48 et 55, les mots insérés à la main sont donnés
par \add{} et les mots tapés biffés par \struck{} ; les soulignements du
tapuscrit sont rendus par l'italique. La main qui annote paraît être celle
de la lettre des pages 8 à 15 ; la page ne le dit pas. En haut à droite,
d'une autre main, droite : « 9/74 Envoyé par Grothendieck »}

\textbf{Esquisse d'une théorie des Gr-Catégories}

par Mme Hoang Xuan Sinh

Nous donnons un résumé de quelques résultats sur les (Gr)-catégories,
faisant l'objet d'un travail détaillé que l'auteur compte présenter
prochainement comme thèse de doctorat [\ ].

\subsection*{1. Structure des (Gr)-catégories.}

Notre terminologie est celle de Saavedra [\ ]. Nous nous intéressons à des
catégories $C$ munies d'une opération binaire $(X,Y) \mapsto X \otimes Y$
(foncteur de $C \times C$ dans $C$) associative et unitaire à isomorphisme
donné près (satisfaisant des conditions de compatibilité explicitées dans
loc.\ cit.), appelées aussi $\otimes$-catégories AU
(associatives-unitaires). On dit qu'une telle catégorie est une
(Gr)-catégorie si \add{c'est un groupoïde, et si} tout objet $X$ de $C$ est
« inversible », i.e.\ admet un objet « inverse » $Y = X^{-1}$ (satisfaisant
$X \otimes Y \simeq Y \otimes X \simeq 1_C$, où $1_C$ désigne l'objet unité
de $C$). Les exemples abondent :

\textbf{Exemples.} 1. $X$ étant un espace topologique \add{ponctué par $x
\in X$}, on prend pour $C$ la catégorie des lacets de $X$ en $x$, avec comme
morphismes les classes d'homotopie d'homotopies entre lacets, comme
opération $\otimes$ la composition des lacets (qui n'est pas associative,
mais associative « à homotopie près »). Variante : $E$ est un espace de Hopf
associatif (ou simplement homotopiquement associatif en un sens suffisamment
fort), $C$ la catégorie dont les objets sont les points de $E$, les
morphismes les classes d'homotopies de chemins entre points de $E$, la loi
$\otimes$ étant induite par la loi de composition de $E$. (NB Lorsque $E$
admet un espace classifiant $X$, on retrouve essentiellement l'exemple
précédent.) \struck{\ill{}}

\note{le tapuscrit porte « Exemples 1. », puis « Exemple 2. » et « Exemple
3. » ; le mot « Exemple » est biffé à la main devant 2 et 3, et le
$s$ de « Exemples » l'est aussi : on lit donc une seule liste. Après la
parenthèse de l'exemple 1, une ligne tapée est biffée au point d'être
illisible}

\marginal{cet exemple est un cas particulier du suivant, en prenant
\uncertain{$E = \Omega(X,x)$} (espace des lacets)}

\note{note marginale à gauche de l'exemple 1, une flèche la rattachant à la
ligne « d'homotopies entre lacets » ; la lettre devant $\Omega(X,x)$ est
douteuse, et $\Omega$ porte peut-être un exposant}

2. Si $F$ est un faisceau sur un espace topologique (ou plus généralement
sur un topos), la catégorie $C$ des torseurs sous $F$, munie de la
composition de Baer, est une (Gr)-catégorie (et même une catégorie de
Picard stricte, cf.\ plus bas). On peut considérer la catégorie $C$ des
Modules inversibles sur un espace (ou topos) \add{localement} annelé
\struck{\ill{}} $(X, \underline{O}_X)$ comme le cas particulier
correspondant au cas $F = \underline{O}_X^{*}$.

3. Si $A$ est une catégorie, la sous-catégorie pleine \add{$C$} de
$\underline{\mathrm{Hom}}(A,A)$, formée des équivalences de $A$ avec
elle-même, munie de l'opération de composition des foncteurs, est une
(Gr)-catégorie. Au lieu de prendre pour $A$ une catégorie, on peut \add{plus
généralement} prendre pour $A$ un objet d'une 2-catégorie quelconque. Si
p.ex.\ $F$ est un faisceau en groupes (pas nécessairement commutatif) sur un
espace topologique (ou un topos \struck{$X$}),

\page{2}
prenant pour $A$ la « champ » sur $X$ formé des $F$-torseurs à droite, la
(Gr)-catégorie des auto-équivalences de $A$ avec lui-même s'interprète comme
la catégorie des « bitorseurs » sous $F$, i.e.\ des faisceaux \add{$P$} sur
lesquels $F$ opère à la fois à droite et à gauche, ces opérations commutant
et chacune d'elles faisant de $P$ un torseur (à droite ou à gauche) sous
$F$, --- la composition $\otimes$ étant la composition de Baer évidente.
Lorsque $F$ est encore de la forme $\underline{O}_X^{*}$ ($\underline{O}_X$
étant un faisceau d'anneaux, qu'on ne suppose plus nécessairement
commutatif) ces bitorseurs s'interprètent aussi en termes de bi-Modules
\add{«} inversibles \add{»} sous $\underline{O}_X$.

\emph{Structure.} Soit $C$ une (Gr)-catégorie, on lui associe

\begin{itemize}
\item[a)] le groupe $\pi_0(C)$ des classes d'isomorphisme d'objets de $C$,
\item[b)] le groupe $\pi_1(C)$ des automorphismes de $1_C$ (objet unité de
$C$)
\item[c)] une action de $\pi_0(C)$ sur $\pi_1(C)$, en associant à tout objet
$X$ de $C$ l'automorphisme $\rho(X)$ de $1_C$ déduit des deux isomorphismes
\[ \mathrm{Aut}(1_C) \rightrightarrows \mathrm{Aut}(X) \]
donnés par $u \mapsto u \otimes \mathrm{id}_X$ et $u \mapsto \mathrm{id}_X
\otimes u$.
\end{itemize}

On prouve que $\pi_1(C)$ est un groupe \emph{commutatif} et que l'on obtient
bien par c) une structure de $\pi_0(C)$-module sur celui-ci. Ceci posé, si
on choisit pour tout $a \in \pi_0(C)$ un représentant $L_a \in \mathrm{Ob}\,
C$ de $a$, et pour deux $a, b$ un isomorphisme
\[ \varphi_{a,b} : L_a \otimes L_b \simeq L_{ab} , \]
alors pour trois éléments $a, b, c \in \pi_0(C)$, l'isomorphisme
d'associativité
\[ (L_a \otimes L_b) \otimes L_c \simeq L_a \otimes (L_b \otimes L_c) , \]
compte tenu des isomorphismes $\varphi_{a,b}$, $\varphi_{ab,c}$,
$\varphi_{b,c}$ et $\varphi_{a,bc}$, peut s'interpréter comme un
isomorphisme
\[ L_{abc} \simeq L_{abc} , \]
ou encore comme la tensorisation à gauche avec un élément bien déterminé
\[ f(a,b,c) \in \pi_1(C) . \]
Les $\varphi_{a,b}$ étant choisis, on voit donc que la donnée d'un
isomorphisme d'associativité fonctoriel $(L \otimes L') \otimes L'' \simeq
L \otimes (L' \otimes L'')$ équivaut à la donnée de l'application
\[ f : \pi_0 \times \pi_0 \times \pi_0 \longrightarrow \pi_1 . \]
On vérifie alors que l'axiome standard d'autocompatibilité d'une donnée
d'associativité \add{(axiome du pentagone)} s'exprime précisément par la
condition que $f$ soit un \emph{3-cocycle} du groupe $\pi_0$ à valeurs dans
le groupe $\pi_1$. D'autre part, l'indétermination dans le choix du système
d'isomorphismes $\varphi_{a,b}$ est précisément donnée par une 2-cochaine
arbitraire, et on voit que si on change $\varphi$ par une 2-cochaine $g$,
$f$ est changé en $f + dg$ --- donc l'ensemble des $f$ correspondants à des
choix différents de $\varphi$ est exactement une classe de 3-cohomologie

\page{3}
\[ k(C) \in H^3(\pi_0(C), \pi_1(C)) . \]
En précisant ces réflexions, on trouve que la classification, à
$\otimes$-équivalence près, des (Gr)-catégories, se fait précisément en
termes d'un groupe $\pi_0$, d'un groupe commutatif $\pi_1$ sur lequel
$\pi_0$ opère, et d'une classe de cohomologie (qui peut être prise
arbitraire) dans $H^3(\pi_0, \pi_1)$. La loi de groupe du $H^3$ admet
d'ailleurs une interprétation « géométrique » à la Baer, en termes
d'opérations sur les (Gr)-catégories.

\emph{Cas particuliers} Dans l'exemple 1, on \add{doit} retrouve\supplied{r}
bien sûr le premier invariant de Postnikoff $\in H^3(\pi_1(X), \pi_2(X))$,
avec une interprétation « géométrique » nouvelle. Dans l'exemple 2, on
trouve $\pi_0 = H^1(X,F)$, $\pi_1 = H^0(X,F)$ et l'élément du $H^3$ est nul
(du fait qu'on a une « catégorie de Picard stricte », cf.\ \add{Deligne}
[\ ]). Dans le dernier exemple, on a $\pi_0$ = groupe des classes
d'isomorphisme d'autoéquivalences de $A$, $\pi_1$ = groupe des
automorphismes du foncteur identique de $A$, et l'invariant $k(C)$ semble
nouveau (même dans le cas où $C$ est la catégorie des biModules inversibles
sous un Anneau $\underline{O}_X$, où $\pi_0$ peut donc s'interpréter comme
un « groupoide de Brandt » et $\pi_1$ comme le groupe des « unités
centrales » de $\underline{O}_X$). Signalons cependant que lorsque $A$ est
la catégorie des torseurs sous un groupe (ordinaire) $F$, on trouve $\pi_0$
= groupe des automorphismes extérieurs de $F$, $\pi_1$ = centre de $F$, et
l'élément $k(C)$ du $H^3$ n'est autre que le classique invariant de
Mac-Lane. Il resterait à \struck{\ill{}} \add{développer} un lien
\emph{direct} \add{géométrique} entre la théorie de Mac-Lane des extensions
de groupes non commutatifs et des liens (« kernels » dans sa terminologie)
auxquels ils sont associés, et la théorie \struck{\ill{}} esquissée ici des
(Gr)-catégories et de leurs splittages.

\note{« développer » est écrit au-dessus d'un mot tapé biffé (« établir »,
sans certitude), et « géométrique » au-dessus de « direct » souligné, sans
que celui-ci soit biffé : on ne sait si l'un remplace l'autre}

\subsection*{2. Catégories de Picard.}

On appelle « catégorie de Picard » une $\otimes$-catégorie ACU
(associative-commutative-unitaire, cf.\ [\ ]) qui est une (Gr)-catégorie
pour sa $\otimes$-structure AU. On dit que c'est une catégorie de Picard
\emph{stricte} si pour tout $X \in \mathrm{Ob}\, C$, la symétrie de $X
\otimes X$ est l'identité. Ces dernières catégories sont étudiées
systématiquement (dans le contexte plus général des « champs de Picards »)
par Deligne [\ ].

Soit $C$ une catégorie de Picard. En tant que (Gr)-catégorie, $C$ possède
des invariants $\pi_0$ et $\pi_1$, mais dans le cas actuel $\pi_0$ est aussi
commutatif et agit trivialement sur $\pi_1$. On fera attention ici qu'une
équivalence de (Gr)-catégories entre deux catégories de Picard $C, C'$ n'est
pas nécessairement compatible avec les contraintes de commutativité, donc
n'est pas nécessairement une équivalence de catégories de

\page{4}
Picard --- on peut dire que la classification « au sens Picard » est plus
fine que la classification au sens des (Gr)-catégories.

Associant à tout $X \in \mathrm{Ob}\, C$ la symétrie de $X \otimes X$,
interprété comme étant la tensorisation par un élément bien déterminé $s(X)$
de $\pi_1(C) = \mathrm{Aut}(1_C)$, on trouve un \struck{application}
\add{homomorphisme} canonique
\[ s(C) = s : \pi_0 \longrightarrow {}_2(\pi_1) , \]
où ${}_2(\pi_1)$ désigne le sous-groupe des éléments $x$ de $\pi_1$ tels que
$2x = 0$. Le résultat principal de la théorie de structure des catégories de
Picard peut s'énoncer en disant que $C$ est déterminée (à équivalence
\add{non canonique} de catégories de Picard près) par la donnée des groupes
commutatifs $\pi_0$ et $\pi_1$, et l'homomorphisme $s : \pi_0 \to
{}_2(\pi_1)$. Dire que cet homomorphisme est nul signifie d'ailleurs que $C$
est une catégorie de Picard \emph{stricte}, et on retrouve le résultat de
[\ ] suivant lequel une catégorie de Picard stricte est déterminée (à
équivalence près) par la connaissance de ses seuls invariants $\pi_0$ et
$\pi_1$.

\note{les indices des $\pi$ dans la formule $s : \pi_0 \to {}_2(\pi_1)$
sont retouchés à la main sur le tapuscrit, qui les avait tapés autrement ;
on donne la formule telle que la ligne suivante la fixe}

\emph{Remarques.} La théorie de loc.\ cit.\ nous permet, plus précisément,
d'interpréter la 2-catégorie des catégories de Picard strictes en termes de
la sous-catégorie pleine de la catégorie dérivée de (Ab) (cat.\ des groupes
abéliens), formée des complexes $K_{\bullet}$ tels que $H_i(K_{\bullet}) =
0$ pour $i \neq 0, 1$. Le résultat précédent, assez grossier, se réduit à
la remarque qu'un objet de cette catégorie est déterminé à isomorphisme (non
unique !) près par la connaissance de son $H_0$ ($= \pi_0$) et $H_1$ ($=
\pi_1$). Il y aurait lieu de développer également pour les catégories de
Picard non nécessairement strictes, et pour les (Gr)-catégories, une théorie
de structure pour la 2-catégorie qu'ils forment. De plus, ces
\struck{résultats} \add{développements} encore dans les limbes, tout comme
ceux que \struck{\ill{}} nous venons d'esquisser, devraient être un jour
\struck{\ill{}} étudiés dans le cadre des champs (en (Gr)-catégories, ou en
catégories de Picard) tout comme dans \add{Deligne} [\ ].

\note{les $H$ de cette remarque, $H_i(K_{\bullet})$, $H_0$, $H_1$, portent
chacun une marque à la main sur l'indice, dont l'intention n'est pas
établie --- peut-être un passage de l'indice à l'exposant ; on les donne
tels que tapés, comme la page 5 les tape encore. Un « alors » manuscrit
court sous la formule $H_i(K_{\bullet}) = 0$, sans signe d'insertion
lisible}

\subsection*{3. Catégories de Picard enveloppantes.}

Soient $C$ une $\otimes$-catégorie AUC (par exemple une catégorie avec
objet final et admettant des produits cartésiens de deux éléments), $S$ une
partie de $\mathrm{Ob}\, C$. On se pose le problème universel d'envoyer (par
un $\otimes$-foncteur AUC) $C$ dans une \supplied{$\otimes$-}catégorie \add{AUC}
\struck{\ill{}} \add{$G$}, de telle façon que pour tout $X \in S$, $F(X)$
soit un objet \emph{inversible} de $G$. Le résultat principal ici est que ce
problème admet toujours une solution, qu'on pourra appeler la
$\otimes$-catégorie AUC « de fractions » par rapport à l'ensemble $S$, et
noter éventuellement $S^{-1}C$. On trouve en fait que pour toute $G$, le
foncteur naturel

\page{5}
\[ \underline{\mathrm{Hom}}_{\otimes \mathrm{AUC}}(S^{-1}C, G)
   \longrightarrow \underline{\mathrm{Hom}}^{S}_{\otimes \mathrm{AUC}}(C, G) \]
(où l'exposant $S$ dans le deuxième membre dénote la sous-catégorie pleine
formée des $\otimes$-foncteurs AUC $F : C \to G$ tels que $F(X)$ soit
inversible pour tout $X \in \mathrm{Ob}\, S$) est un \emph{isomorphisme de
catégories}.

Nous n'explicitons pas ici la construction, assez naturelle, mais dont
l'explicitation pas à pas, \struck{et} \add{avec} la vérification de toutes
les compatibilités voulues, est fort longue et fastidieuse. Signalons
seulement ici que lorsque $S$ est l'ensemble de tous les objets de $C$,
$S^{-1}C$ est elle-même une catégorie de Picard (pas nécessairement
stricte), qu'on pourra appeler la « catégorie de Picard enveloppante » de
$C$. C'est sans doute le cas le plus important. Les invariants $\pi_0$ et
$\pi_1$ de cette catégorie enveloppante méritent la notation $K^0(C)$ et
$K^1(C)$ respectivement. Par exemple, lorsque $C$ est la catégorie des
modules projectifs de type fini sur un anneau $A$, \add{avec l'opération de
somme,} on peut démontrer que ces groupes ne sont autres que les classiques
invariants $K^0(A)$ et $K^1(A)$ étudiés par Grothendieck et par
Whitehead-Bass. Le même résultat est sans doute valable pour toute catégorie
additive $C$ (munie de l'opération somme). Cela donne une interprétation
remarquable des invariants $K^0$ et $K^1$ en termes d'une « catégorie
stabilisée » (de la catégorie additive $C$), laquelle est une catégorie de
Picard. Signalons seulement que, cette catégorie stabilisée (même dans le
cas particulier des modules projectifs de type fini sur l'anneau $A$)
n'étant pratiquement jamais une catégorie de Picard \emph{stricte}, ces
considérations semblent rendre peu probable qu'on arrive à construire un
objet canonique dans la catégorie dérivée de la catégorie (Ab), qui soit un
complexe de chaines satisfaisant $H_0(K_{\bullet}) = K^0$,
$H_1(K_{\bullet}) = K^1$ --- puisqu'une telle construction (grâce à la
théorie de Deligne) nous fournirait précisément une catégorie de Picard
\emph{stricte} canonique dont les invariants seraient $K^0$ et $K^1$.

On peut difficilement échapper à la question d'une interpétation
géométrique analogue des invariants $K^i$ supérieurs. Il y a tout lieu de
croire qu'une telle interprétation est à chercher dans la théorie des
$n$-catégories de Picard (en un sens convenable) et la construction des
$n$-catégories de Picard enveloppantes d'une $\otimes$-catégorie AUC (ou
d'une $\otimes$-$n$-catégorie AUC) arbitraire $C$, dont les invariants
« homotopiques » $\pi_i$ seraient les $K^i$ cherchés. De tels
développements obligerai-

\struck{Comme autre direction de recherche ouverte, suggérée par les
résultats esquissés ici, signalons la construction d'une
$\otimes$-catégorie de fractions d'une $\otimes$-catégorie AU (sans donnée
de commutativité néces-}

\note{ce paragraphe, qui s'arrête en bas de page au milieu d'un mot, est
biffé de grands traits obliques ; la page 6 enchaîne sur la phrase qui le
précède}

\page{6}
ent sans doute à tirer au clair les relations très étroites qu'on pressent
entre les notions de $n$-catégorie (et de $\infty$-catégorie) d'une part,
celle d'ensemble simplicial (ou d'espace topologique …) d'autre part, tout
comme la relation entre $n$-Gr-catégories (et $\infty$-Gr-catégories) et
groupes simpliciaux (ou groupes topologiques), enfin entre $n$-catégories de
Picard \emph{strictes} et complexes de chaines tronqués à la dimension $n$.
Il y aurait à s'attendre à une fusion plus ou moins complète entre ces trois
visions : algèbre homologique, algèbre homotopique, algèbre catégorique,
dans une vision commune (qui pourrait bien prendre le nom d'« algèbre
homologique non commutative », étant entendu que ce qui porte actuellement
ce nom n'est que l'amorce de ce qui \struck{reste à} \add{doit} venir …)

Comme dernière question, plus terre à terre, signalons celle liée à la
remarque \struck{que} \add{suivante. Lorsque} \struck{\ill{}} $C$ est la
catégorie des modules projectifs de type fini sur un anneau $A$
\emph{commutatif} (ou plus généralement une catégorie additive munie d'une
opération $\otimes$ \struck{\ill{}} \add{ACU} \emph{biadditive}), alors la
catégorie de Picard enveloppante, en plus de sa loi de composition évidente
de « somme » (notée ici $\oplus$) donnée \struck{\ill{}} \add{par} sa
définition même, est munie d'une opération $\otimes$ provenant de
l'opération de produit tensoriel \add{ordinaire} des modules, qui en fait
également une $\otimes$-catégorie ACU (mais évidemment pas une
(Gr)-catégorie), liée par des isomorphismes de distributivité évidents
\struck{aux} \supplied{à l'}opération $\oplus$. Il y aurait lieu d'étudier
systématiquement les catégories $A$ ainsi munies de deux opérations $\oplus$
et $\otimes$, toutes deux AUC (des données supplémentaires), avec une donnée
de distributivité
\[ X \otimes (Y \oplus Z) \simeq X \otimes Y \oplus X \otimes Z \]
et des compatibilités convenables pour celle-ci ; et de tenter de faire une
théorie de structure pour de telles catégories qui sont de Picard pour la
structure $\oplus$. Cela implique en particulier que $\pi_0(A)$ est un
anneau unitaire commutatif, et $\pi_1(A)$ un module sous ce dernier. Quels
autres invariants faut-il introduire pour caractériser $A$ à équivalence
près respectant toutes les structures ?

\section*{A.\ Grothendieck à P.\ Deligne, de Lodève, 27.8.74 (pages 8 à 15)}

\page{8}
\note{Les pages 8 à 11 sont la lettre que le dossier 134-1 porte en ses
pages 3 à 6 ; cet exemplaire-ci se poursuit quatre pages de plus. La date
est pâle : « 7.8.74 » se lit nettement, mais devant le 7 ne reste qu'un
court trait, le pied d'un 2 effacé ; l'exemplaire du dossier 134-1 porte
lisiblement « 27.8.74 »}

\marginal{Je n'ai pas retrouvé la lettre sur les Gr-champs que tu cites}

\note{note écrite verticalement dans la marge de gauche, d'une main que rien
ne distingue de celle de la lettre ; la page ne dit pas si elle est de lui
ou du destinataire}

Lodève le \uncertain{2}7.8.74

Cher Deligne,

Étant peut-être empêché \uncertain{par une jambe} d'assurer un cours de 1er
cycle au 1er trimestre, je vais peut-être à la place faire un petit
séminaire d'algèbre, et envisage de le faire sur les fourbis de Mme Sinh,
éventuellement transposés dans le contexte des « champs ». À ce propos, je
tombe sur le truc suivant, qui pour l'instant reste heuristique. Si $M$,
$N$ sont deux faisceaux abéliens sur un topos $X$, et
\[ \tau_{\leq 2}\, R\underline{\mathrm{Hom}}(M,N) = E(M,N) \]
est le complexe ayant les invariants
\[ \begin{cases}
\underline{H}^i = \underline{\mathrm{Ext}}^i(M,N) & \text{pour } 0 \leq i \leq 2 \\
\underline{H}^i = 0 & \text{si } i \notin [0,2]
\end{cases} \]
il doit y avoir un triangle distingué canonique $(T)$ :

\begin{tikzcd}
& \underline{\mathrm{Hom}}(M, {}_2 N)[-2] \arrow[dl] & \\
E(M,N) \arrow[rr] & & E'(M,N) \arrow[ul]
\end{tikzcd}

\page{9}
donc $E'(M,N)$ est un complexe dont les invariants $\underline{H}^i$ sont
ceux de $E(M,N)$ en degré $i+2$, et qui en degré 2 donne lieu à une suite
exacte
\[ (*) \qquad 0 \longrightarrow \underline{\mathrm{Ext}}^2(M,N)
   \longrightarrow \underbrace{\underline{H}^2(E'(M,N))}_{P(M,N)}
   \xrightarrow{\ \sigma\ } \underline{\mathrm{Hom}}(M, {}_2 N)
   \longrightarrow 0 \]
Heuristiquement, $E'(M,N)$ est le complexe qui exprime le « 2-champ de
Picard "strict" » formé des 1-champs de Picard (pas nécessairement stricts)
« épinglés » par $M$, $N$ sur des objets variables de $X$, en admettant que
la théorie pour les 1-champs de Picard stricts s'étend aux 2-champs de
Picard stricts (ce qui pour moi ne fait guère de doute) ; de même $E(M,N)$
correspond aux champs de Picard \emph{stricts} épinglés par $M$, $N$.

La suite exacte $(*)$ se construit en tous cas canoniquement « à la
main »), où le terme médian est le faisceau des \uncertain{classes} à
« équivalence » près des

\note{« épinglés par $M$, $N$ » et « stricts » (après « 1-champs de
Picard ») sont insérés en interligne. La parenthèse fermante après « à la
main » est sur la page, sans ouvrante}

\page{10}
\struck{catégories} \add{champs} de Picard épinglés par $M$, $N$, $\sigma$
étant l'invariant qui s'obtient en associant à toute section $L$ d'un champ
de Picard la symétrie de \struck{\ill{}} $L \otimes L$, interprétée comme
section de ${}_2 N$. Je sais prouver (sauf erreur) que tout hom.\ $M \to
{}_2 N$ provient d'un champ de Picard convenable (épinglé par $M$, $N$)
(a priori l'obstruction est dans $\mathrm{Ext}^3(X; M, N)$, mais un
argument « universel » prouve qu'elle est nulle). Cela prouve que
l'extension $(*)$ est bien proche d'être splittée : toute section du
troisième faisceau, sur un objet quelconque de $X$, se relève --- en
d'autres termes l'extension a une section « ensembliste ». Bien sûr, il y a
mieux en fait : toute section sur un $U \in \mathrm{Ob}\, X$ « provient »
d'un élément de $\mathbb{H}^2(U, E'(M,N))$ (hypercohomologie
$\mathbb{H}^2$).

\note{« champs » est écrit au-dessus de « catégories » biffé ; « s'obtient
en » et « $L$ » sont insérés en interligne. L'exemplaire du dossier 134-1
porte la même insertion. Devant
$L \otimes L$ un symbole est biffé}

\page{11}
\emph{Exemple} Soit $A$ un anneau sur $X$, soient $M$, $N$ respectivement
les faisceaux $K^0$ et $K^1$ associés au champ additif des $A$-Modules
projectifs de t.f.\ (p.ex.). Alors le raisonnement de Mme Sinh nous fournit
un champ de Picard épinglé par $M$, $N$, d'où une section \struck{de}
canonique du terme médian $P(M,N)$ de $(*)$.

\emph{NB} Tout ce qui précède a les fonctorialités évidentes en $M$, $N$,
$X$ …

\emph{Question} : Le triangle exact $(T)$ et la suite exacte $(*)$ sont
ils connus par les \uncertain{compétences} (Quillen, Breen, Illusie …) ?
Connaissent-ils des variantes « supérieures » ? (Un principe
« géométrique » pour les obtenir pourrait être via des $n$-champs de Picard
non néc.\ stricts …)

\page{12}
\note{L'ordre des archivistes n'est probablement pas l'ordre d'écriture :
cette page renvoie à « $B_e \to B_G$ plus haut », et $B_G$ n'apparaît qu'aux
pages 14 et 15 ; la page 13 s'achève sur « J'ai » et la page 14 s'ouvre sur
« Je te signale ». La suite 8--11, 14--15, 12--13 se lit sans rupture, et
le feuillet qui suivait la page 13 n'est pas dans le dossier. On transcrit
dans l'ordre des archivistes}

Je profite de l'occasion pour soulever une question sur la « cohomologie
relative ». Soit $q : X \to Y$ un morphisme de topos. Si $F$ est un faisceau
abélien (ou un complexe d'iceux) sur $Y$, peut-on définir
\emph{fonctoriellement} en $F$ la cohomologie relative $R\Gamma(Y \bmod X,
F)$ (de la catégorie dérivée de $\mathrm{Ab}(Y)$ dans celle de
$\mathrm{Ab}$) ? L'interprétation « géométrique » en termes d'opérations
sur des $n$-champs de Picard ($n$ « grand ») suggère que ça doit exister.
Mais je ne vois de constructions évidentes « à la main » que dans les deux
cas extrêmes :

a) $q$ est « $(-1)$-acyclique » i.e.\ pour tt $F$ sur $Y$, $F \to q_* q^* F$
est injectif (NB c'est le cas de $Y_{/P} \to Y$ où $P \to e_Y$ est un
épimorphisme --- c'est donc le cas de $B_e \to B_G$ \uncertain{plus haut}).
On prend
\[ R\Gamma_Y\bigl(\mathrm{Coker}(F \to q_*(C(q^*(F))))[-1]\bigr) \]

\note{« résol.\ inj.\ » est écrit sous le $C$, qui désigne donc une
résolution injective ; l'indice $Y$ de $R\Gamma_Y$ est empâté}

\page{13}
b) $\forall F$ injectif sur $Y$, $q^*(F)$ est injectif et $F \to q_* q^* F$
est un épimorphisme (ex : $q$ inclusion d'un ouvert $U \subset e_Y$). On
prend
\[ R\Gamma_Y\bigl(\mathrm{Ker}(C(F) \to q_* q^*(C(F)))\bigr) \]

\note{« rés.\ inj.\ » sous le premier $C$ ; après $q_* q^*$ le trait
d'ouverture est fait comme un $C$, et l'on pourrait lire $q_* q^* C(C(F))$}

Dans le cas général, la difficulté provient du fait que le cône d'un
morphisme de complexes (tel que
$F \to q_*(q^*(F))$ \struck{\ill{}} )
n'est pas fonctoriel (dans la catégorie dérivée) p.r.\ à la flèche dont on
veut prendre le cône.

Et pourtant, dans les cas particuliers actuels, il devrait y avoir un choix
fonctoriel. Est-ce évident ?

\emph{Question pour Illusie} : Dans sa théorie des déformations de schémas
en groupes plats, il tombe sur des $H^3(B_G/X,\ )$ resp.\ des
$\mathrm{Ext}^2(X,\ )$. Peut-on court-circuiter sa théorie via la théorie
(supposée écrite) des Gr-champs --- resp.\ via la théorie des champs de
Picard ? J'ai

\note{les seconds arguments des $H^3$ et du $\mathrm{Ext}^2$ sont laissés
en blanc sur la page ; la lettre dans le $\mathrm{Ext}^2$ est douteuse
($X$ ou $K$, avec un point souscrit). Le « J'ai » final n'a pas de suite
dans le dossier}

\page{14}
Je te signale que j'ai réfléchi aux Gr-champs sur $X$. Si $G$ est un Groupe
sur $X$, $N$ un $G$-Module les Gr-champs sur $X$ « épinglés par $G$, $N$ »
forment a priori une 2-catégorie et même une 2-catégorie de Picard stricte,
grâce à l'opération évidente à la Baer. On trouve que le complexe \add{(de
cochaînes)} tronqué à l'échelon 2 qui leur correspond est le tronqué
\[ \tau_{\leq 2}\bigl( R\Gamma(B_G \bmod X, N)[1] \bigr) \]
(\emph{NB} la cohomologie de $R\Gamma(B_G \bmod X, N)$ commence en degré 1)
Plus géométriquement, un Gr-champ sur $X$ épinglé \add{par $(G,N)$} est
essentiellement « la même chose » qu'une 2-gerbe sur $B_G$, liée par $N$,
et \struck{\ill{}} munie d'une trivialisation au dessus de $X \approx B_e =
(B_G)_{/P}$ (où $P$ est le « objet de $B_G$ "torseur universel sous $G$" »).
Ces 2-gerbes forment en fait une 3-catégorie de Picard a priori, mais il se

\note{« (de cochaînes) » est inséré en interligne, « par $(G,N)$ » dans la
marge de gauche avec un trait de renvoi}

\page{15}
trouve que dans celle-ci, les 3-flèches sont triviales (i.e.\ si source =
but, ce sont des identités) --- cela \uncertain{un} fait qui exprime
$H^0(B_G/X, N) = 0$ (i.e.\ $H^0(B_G, N) \to H^0(X, N)$ injectif …). Donc la
3-catégorie peut être regardée comme une 2-catégorie --- et « c'est » celle
des Gr-champs sur $X$ épinglés par $G$, $N$. Si on veut localiser sur $X$,
et décrire le 2-champ \add{de Picard} sur $X$ des champs de Picard (sur des
objets variables de $X$) épinglés par $G$, $N$, on trouve qu'il est exprimé
par le complexe
\[ \tau_{\leq 2}\bigl( R p_{G*}\, \mathrm{Coker}(N \to R q_{G*}\, C(q_G^* N)) \bigr) \]
où $p_G : B_G \to X$ et $q_G : B_e \simeq X \simeq (B_G)_{/P} \to B_G$.
Toutes ces descriptions étant compatibles avec des variations de $G$, $N$,
$X$, cela donne en principe une description de la 2-catégorie des
Gr-champs, avec $X$, $G$, $N$ variables …

\note{« rés.\ inj.\ » au-dessus du $C$ de la formule}

\section*{A.\ Grothendieck à Deligne, Giraud et Verdier, de Villecun, 23.6.1974 (pages 16 et 17)}

\page{16}
\note{Lettre tapée à la machine, signée à la main ; une insertion
manuscrite au dernier paragraphe}

Villecun le 23.6.1974

Chers Deligne, Giraud, Verdier,

Vous savez peut-être qu'une mathématicienne vietnamienne, Hoang Xuan Sinh,
a travaillé depuis mon séjour au Vietnam en 1967 à un sujet de recherche que
je lui avais proposé, savoir une théorie de structure des (Gr)-catégories,
et la construction d'une (Gr)-catégorie (ou plutôt catégorie de Picard)
enveloppante pour une catégorie avec produit $\otimes$ associatif et
commutatif. J'ai reçu récemment une copie du manuscript en forme (dont la
confection s'est faite au prix de mille difficultés --- il est douteux qu'il
soit possible d'avoir d'autres copies dans un avenir proche, à moins d'en
faire d'après cet exemplaire-là), et j'ai envoyé un rapport officiel au
sujet de ce travail à M.\ Ta Quang Buu, ministre de l'enseignement
supérieur et technique de la RDV, rapport dont je vous envoie une copie
ci-joint. Mme Sinh, avec mon accord, a demandé à présenter ce travail comme
thèse de doctorat, et à faire sa soutenance en France. À l'heure où
j'écris, je n'ai pas connaissance d'un accord des autorités vietnamiennes
pour l'un ni pour l'autre ; il y a eu un certain nombre de difficultés
soulevées, qu'il est inutile de détailler ici, et qui j'espère vont
s'aplanir. J'ai appris par ailleurs que Mme Hoang Xuan Sinh figurerait sur
la liste officielle de la délégation vietnamienne au congrès international
de mathématiciens à Vancouver du 20 août 1974. Si c'est bien ainsi, il
serait très souhaitable qu'elle puisse profiter de cette occasion pour
faire sa soutenance à Orsay au retour de Vancouver, fin Septembre ou début
Octobre, et il me semble plausible que les autorités vietnamiennes y
donneront leur accord. Aussi il serait assez urgent, vu les lenteurs
administratives, de constituer un jury. Par la présente, je vous demande si
vous seriez disposé en principe à faire partie (avec moi) de ce jury de
thèse --- ce qui suppose évidemment que vous vous trouviez dans la région
parisienne au moment prévu pour la soutenance. À celui qui m'en ferait la
demande, je peux envoyer le texte de la thèse complète, ou un résumé assez
détaillé rédigé par Mme Hoang Xuan Sinh, ou enfin un résumé moins détaillé
de 13 pages manuscrites (rédigées en vue d'une publication aux CR, mais
beaucoup trop long à cet effet). Ce serait gentil de me répondre rapidement,
et en cas d'impossibilité, de voir autour de vous s'il y a des gens « dans
le bain » \struck{\ill{}} \add{de choses} que fait Mme Sinh, et qui seraient
intéressés à faire partie du jury.

\note{le rapport annoncé « ci-joint » n'est pas dans le dossier. Deux mots
tapés sont biffés devant « que fait Mme Sinh », et « de choses » écrit
au-dessus}

\page{17}
Pour ce qui est des formalités administratives, c'est le frère de Mme
\uncertain{Hoan Xuan Man}, qui habite à Antony, qui s'en occupera pour elle.
A toutes fin utile, je vous passe son adresse :

\note{le nom est tapé ainsi, sic. La ligne suivante, l'adresse postale et le
numéro de téléphone du frère, n'est pas reproduite}

Dans l'attente d'une réponse prochaine, bien cordialement

Schurik

\note{la signature est manuscrite}

PS N'ayant pas l'adresse de Giraud, je demande à Verdier s'il peut bien lui
transmettre la lettre et le rapport. Je pense que celui-ci doit pouvoir
servir comme rapport de thèse aussi vis à vis de l'administration
universitaire en France.

\end{document}
